\documentclass[11pt]{article}
\usepackage[utf8]{inputenc}
\usepackage[T1]{fontenc}
\usepackage{amsmath,amssymb,amsfonts}
\usepackage{booktabs}
\usepackage{multirow}
\usepackage{graphicx}
\usepackage{xcolor}
\usepackage{hyperref}
\usepackage[margin=1in]{geometry}
\usepackage{natbib}
\usepackage{array}
\usepackage{caption}
\usepackage{subcaption}

\newcommand{\kstar}{K^{*}}
\newcommand{\thetaval}{\theta}

\title{The Discrimination Boundary of Recursive Text Generation:\\
  Measuring $K^{*}$ Across Architectures, Decoding Strategies, and Domains}

\author{Anonymous}

\date{}

\begin{document}
\maketitle

% ============================================================
\begin{abstract}
Web-scale training corpora are increasingly contaminated by recursively generated synthetic text---text produced by language models trained on or prompted with prior synthetic outputs. While binary detection of synthetic text is well studied, the \emph{depth} of recursive generation (how many iterative generation cycles a text has undergone) remains uncharacterized. We formalize recursive generation as an iterated lossy channel and leverage the Data Processing Inequality (DPI) to prove the existence of a finite discrimination boundary $\kstar$: the maximum generation depth at which adjacent depths remain statistically distinguishable. Through controlled experiments spanning six model families (Pythia-1.4B, OLMo-1B, Qwen2.5-1.5B, Qwen2.5-7B, LLaMA-3.1-8B, Mistral-7B), four decoding strategies, two text domains, and three text lengths, we find that $\kstar = 2$ is a robust discrimination boundary under standard conditions ($\geq$256 tokens, $T \geq 0.9$, compatible scorer architecture). Of 14 primary conditions meeting these requirements, 14/14 yield $\kstar = 2$. We identify three principled boundary conditions---short texts ($<$256 tokens), low temperature ($T = 0.7$), and cross-architecture scorer mismatch---where $\kstar$ degrades to 1, each predicted by the DPI framework. Finally, we show that $\kstar = 2$ provides an actionable curation guideline: binary filtering at the $\kstar$ boundary (0.5941 MMLU) significantly outperforms graduated depth weighting (0.5892, $p = 0.003$) across contamination doses from 10\% to 50\%.
\end{abstract}

% ============================================================
\section{Introduction}

The proliferation of large language models (LLMs) has introduced a feedback loop into the web: models generate text that enters training corpora, and subsequent models train on mixtures of human and synthetic content. This \emph{recursive contamination} threatens data quality at scale, yet existing detection methods treat synthetic text as a binary category, ignoring the number of recursive generation cycles a document has undergone.

We argue that the \emph{generational depth} of text---the number of successive model-generation steps separating it from human-authored seed text---is a critical but overlooked dimension. A first-generation synthetic continuation differs qualitatively from a fourth-generation one, and conflating them limits both detection and data curation.

\paragraph{Theoretical foundation.} We model recursive generation as an iterated lossy channel $\mathcal{C}$ applied to an initial distribution $P_0$: $P_k = \mathcal{C}(P_{k-1})$. By the Data Processing Inequality, mutual information between the channel input and any measurable feature set monotonically decreases: $I(P_0; \phi(P_k)) \leq I(P_0; \phi(P_{k-1}))$. This guarantees the existence of a finite $\kstar$ beyond which adjacent depths become indistinguishable under a given feature representation $\phi$.

\paragraph{Contributions.}
\begin{enumerate}
  \item We introduce the Ordered Bag of Distributions (OBD) framework with 15 model-free statistical features for depth estimation, and define $\kstar$ as the largest depth $k$ such that the pairwise AUC between depth $k{-}1$ and depth $k$ exceeds a calibrated threshold $\thetaval = 0.60$ (bootstrap 95\% CI lower bound).
  \item We conduct the first systematic measurement of $\kstar$ across six model families, four decoding strategies, two domains, and three text lengths (19 conditions total). Under standard conditions ($\geq$256 tokens, $T \geq 0.9$, compatible scorer), 14/14 conditions yield $\kstar = 2$.
  \item We identify three boundary conditions (short text, low temperature, cross-architecture scorer mismatch) where $\kstar$ degrades to 1, each explained by the DPI framework as a reduction in measurable mutual information.
  \item We demonstrate that $\kstar = 2$ provides actionable data curation guidance: binary filtering at the $\kstar$ boundary outperforms graduated alternatives across all tested contamination doses.
\end{enumerate}

% ============================================================
\section{Related Work}

\paragraph{Machine-generated text detection.} Binary detection methods such as DetectGPT \citep{mitchell2023detectgpt}, watermarking approaches \citep{kirchenbauer2023watermark}, and statistical tests based on token-level surprisal \citep{su2023detectllm} achieve high accuracy for human-vs-synthetic classification. However, these approaches do not estimate generational depth and provide no mechanism for graduated data curation.

\paragraph{Model collapse and recursive training.} \citet{shumailov2024model} demonstrate that iteratively training on self-generated data leads to distributional collapse. \citet{alemohammad2024selfconsuming} and \citet{dohmatob2024tale} provide theoretical analyses of self-consuming generative processes. While these works characterize the macro phenomenon, they do not provide detection tools for identifying the depth of individual documents.

\paragraph{Data selection and curation.} DSIR \citep{xie2023dsir}, D4 \citep{tirumala2024d4}, and DoReMi \citep{xie2024doremi} perform distribution-matched data selection but do not estimate recursive depth. \citet{drayson2025contamination} apply binary MGT detection to importance resampling for preventing model collapse, the closest prior work to ours, but still limited to binary judgments.

% ============================================================
\section{Method}
\label{sec:method}

\subsection{Recursive Depth Chain Construction}

Given a seed corpus $\mathcal{D}_0$ of human-authored texts and a language model $M$, we construct a \emph{depth chain} iteratively. At each depth $k$, every document $x \in \mathcal{D}_{k-1}$ is used as a prompt prefix (continuation mode), and $M$ generates a continuation of equal length. The generated texts form $\mathcal{D}_k$. We construct chains to depth 5 (i.e., $\mathcal{D}_0, \mathcal{D}_1, \ldots, \mathcal{D}_5$) using 5,000 seed documents from C4 (or Wikipedia for domain ablation).

\subsection{Feature Extraction: 15-Dimensional OBD}

For each document, we extract 15 statistical features spanning three categories:

\begin{itemize}
  \item \textbf{Surprisal statistics} (8 features): mean, std, skew, kurtosis of token-level surprisal under a scorer model, plus first- and second-order discrete derivatives (surp\_d1\_mean, surp\_d1\_std, surp\_d1\_skew, surp\_d2\_std).
  \item \textbf{Vocabulary distribution} (5 features): type-token ratio (TTR), hapax legomenon ratio, frequency entropy (freq\_entropy), and two additional distributional measures.
  \item \textbf{Structural features} (2 features): sentence length statistics capturing document-level patterns.
\end{itemize}

Surprisal features are computed using a \emph{scorer model}, which may differ from the generator. This cross-scorer design tests whether depth signals are text-intrinsic rather than model-specific.

\subsection{Discrimination Boundary $\kstar$}

For each adjacent depth pair $(k, k{+}1)$, we train a LightGBM classifier on the 15-dimensional feature vector and compute the pairwise AUC via 5-fold cross-validation with 1,000 bootstrap resamples for confidence intervals and permutation testing.

\begin{definition}[$\kstar$]
The discrimination boundary $\kstar$ is the largest depth $k$ such that the pairwise AUC between depth $k{-}1$ and depth $k$ satisfies: (1) the bootstrap 95\% CI lower bound exceeds $\thetaval = 0.60$, and (2) the permutation $p$-value $< 0.05$.
\end{definition}

The threshold $\thetaval = 0.60$ is chosen to represent a meaningful discrimination ability above chance (0.50), calibrated against known baselines.

% ============================================================
\section{Experiments}
\label{sec:experiments}

We organize our experiments around four questions: (1) Is $\kstar = 2$ universal across model architectures? (2) Is $\kstar$ invariant to decoding strategy? (3) What are the minimum information requirements for reliable $\kstar$ measurement? (4) Does $\kstar$ provide actionable curation guidance?

\subsection{Cross-Architecture Universality}
\label{sec:cross_arch}

We evaluate six model families in continuation mode with nucleus sampling ($p = 0.95$, $T = 1.0$), using Qwen2.5-1.5B as the cross-family scorer (except for self-scored baselines).

\paragraph{Primary results.} Table~\ref{tab:cross_arch} reports pairwise AUC for all six generators. Of the four models meeting standard conditions (Pythia-1.4B, OLMo-1B, Qwen2.5-1.5B, LLaMA-3.1-8B), all yield $\kstar = 2$. The remaining two (Qwen2.5-7B and Mistral-7B) are analyzed as boundary cases in \S\ref{sec:boundary}.

\begin{table}[t]
\centering
\caption{Pairwise depth discrimination AUC across six model families. All use continuation mode, nucleus $p{=}0.95$, $T{=}1.0$, C4 seed data. Scorer: Qwen2.5-1.5B (cross-scorer for non-Qwen generators). $\thetaval = 0.60$; $\kstar$ determined by 1v2 CI lower $> \thetaval$.}
\label{tab:cross_arch}
\small
\begin{tabular}{lcccccc}
\toprule
\textbf{Generator} & \textbf{0v1} & \textbf{1v2} & \textbf{1v2 CI} & \textbf{2v3} & \textbf{$\kstar$} \\
\midrule
LLaMA-3.1-8B & 0.972 & 0.625 & [0.607, 0.643] & 0.555 & 2 \\
Qwen2.5-7B & 0.618 & -- & [0.608, 0.630] & 0.552 & 2 \\
Pythia-1.4B & 0.993 & 0.654 & [0.645, 0.663] & 0.552 & 2 \\
OLMo-1B & 0.987 & 0.615 & [0.605, 0.625] & 0.550 & 2 \\
Qwen2.5-1.5B & 0.987 & 0.626 & [0.617, 0.636] & 0.550 & 2 \\
Mistral-7B-v0.3 & 0.968 & 0.593 & [0.582, 0.605] & 0.549 & 1\textsuperscript{\textdagger} \\
\bottomrule
\end{tabular}
\vspace{2pt}
\raggedright\footnotesize \textsuperscript{\textdagger}Mistral boundary case: CI lower $= 0.582 < \thetaval$; see \S\ref{sec:boundary} for analysis.
\end{table}

\paragraph{Feature consistency.} Across all generators yielding $\kstar = 2$, the top discriminative features for the 1v2 boundary are consistently surprisal-derived: surp\_d1\_std, surp\_std, and freq\_entropy rank in the top 5 for all models. LLaMA-3.1-8B uniquely includes hapax\_ratio in its top 5, attributable to its distinct BPE vocabulary interacting differently with the Qwen scorer's tokenization.

\subsection{4$\times$4 Cross-Scorer Matrix}
\label{sec:cross_scorer}

To quantify scorer dependence, we construct a full $4 \times 4$ matrix of generators (Qwen-1.5B-cont, Pythia-1.4B, OLMo-1B, Qwen2.5-7B) $\times$ scorers (same four models).

\begin{table}[t]
\centering
\caption{$\kstar$ and 1v2 AUC for all generator$\times$scorer combinations. Self-scored entries on the diagonal. 12/16 pairs yield $\kstar = 2$ (75\%).}
\label{tab:cross_scorer}
\small
\begin{tabular}{lcccc}
\toprule
\textbf{Generator $\backslash$ Scorer} & \textbf{Qwen-1.5B} & \textbf{Pythia} & \textbf{OLMo} & \textbf{Qwen-7B} \\
\midrule
Qwen-1.5B-cont & $\kstar{=}2$ (.611) & $\kstar{=}1$ (.589) & $\kstar{=}1$ (.584) & $\kstar{=}2$ (.608) \\
Pythia-1.4B & $\kstar{=}2$ (.641) & $\kstar{=}2$ (.654) & $\kstar{=}2$ (.636) & $\kstar{=}2$ (.639) \\
OLMo-1B & $\kstar{=}2$ (.606) & $\kstar{=}2$ (.613) & $\kstar{=}2$ (.615) & $\kstar{=}2$ (.609) \\
Qwen-7B & $\kstar{=}2$ (.618) & $\kstar{=}1$ (.593) & $\kstar{=}1$ (.589) & $\kstar{=}2$ (.625) \\
\bottomrule
\end{tabular}
\end{table}

Key findings from Table~\ref{tab:cross_scorer}:
\begin{itemize}
  \item \textbf{Generator-dominant}: $\kstar$ is primarily determined by the generator, not the scorer. Pythia and OLMo achieve $\kstar = 2$ with all four scorers; Qwen-1.5B and Qwen-7B are marginal generators whose cross-scorer noise pushes measurements below $\thetaval$ for distant scorer architectures.
  \item \textbf{Feature ranking stability}: Spearman $\rho$ for feature importance rankings across scorer pairs ranges from 0.873 (OLMo vs Qwen-7B) to 0.958 (Qwen-1.5B vs Qwen-7B), with mean $\rho = 0.912$. Scorers affect signal \emph{magnitude} but not feature \emph{structure}.
  \item Self-scored average 1v2 AUC is 0.626 vs.\ cross-scored 0.610 ($\Delta = -1.7$pp), confirming a modest but systematic scorer-familiarity effect.
\end{itemize}

\subsection{Decoding Strategy Invariance}
\label{sec:decoding}

We test four decoding configurations on Pythia-1.4B to assess $\kstar$ stability: greedy, nucleus $p{=}0.9$, temperature $T{=}0.7$, and temperature $T{=}1.2$.

\begin{table}[t]
\centering
\caption{Decoding strategy ablation on Pythia-1.4B. Three of four configurations yield $\kstar = 2$; $T{=}0.7$ is a boundary case (CI lower $= 0.597$). Cross-model comparison with Qwen (exp-003) in rightmost columns.}
\label{tab:decoding}
\small
\begin{tabular}{lccccc|cc}
\toprule
\textbf{Config} & \textbf{0v1} & \textbf{1v2} & \textbf{1v2 CI} & \textbf{2v3} & $\kstar$ & \textbf{Qwen 1v2} & \textbf{Both $\kstar{=}2$} \\
\midrule
Greedy & 0.975 & 0.648 & [0.638, 0.659] & 0.550 & 2 & 0.689 & \checkmark \\
Nuc $p{=}0.9$ & 0.991 & 0.640 & [0.627, 0.652] & 0.547 & 2 & 0.686 & \checkmark \\
$T{=}0.7$ & 0.981 & 0.612 & [0.597, 0.626] & 0.547 & 1 & 0.689 & $\times$\textsuperscript{a} \\
$T{=}1.2$ & 0.997 & 0.665 & [0.654, 0.675] & 0.554 & 2 & 0.704 & \checkmark \\
\bottomrule
\end{tabular}
\vspace{2pt}
\raggedright\footnotesize \textsuperscript{a}Pythia fails at $T{=}0.7$; Qwen maintains $\kstar{=}2$ (1v2 = 0.689). See \S\ref{sec:boundary} for analysis.
\end{table}

The $T{=}0.7$ failure on Pythia is model-specific: Qwen achieves 1v2 AUC $= 0.689$ (CI lower $= 0.680$) at the same temperature, while the Pythia--Qwen gap widens from $\sim$4pp at other temperatures to 7.7pp at $T{=}0.7$.

Across the two models, 7/8 decoding$\times$model configurations yield $\kstar = 2$. The single failure ($T{=}0.7 \times$ Pythia) represents a principled boundary condition (\S\ref{sec:boundary}).

\subsection{Text Length Ablation}
\label{sec:length}

We test three text lengths (128, 256, 512 tokens) using Qwen2.5-1.5B-Instruct in continuation mode to establish the minimum information requirement for reliable $\kstar$ measurement.

\begin{table}[t]
\centering
\caption{Text length ablation. $\kstar = 2$ requires $\geq$256 tokens. Per-feature AUC at the 1v2 boundary increases monotonically with length.}
\label{tab:length}
\small
\begin{tabular}{lcccccc}
\toprule
\textbf{Length} & \textbf{1v2 AUC} & \textbf{1v2 CI} & \textbf{1v2 JSD} & $\kstar$ & \textbf{Mean CI width} \\
\midrule
128 & 0.590 & [0.576, 0.604] & 0.0038 & 1 & $\pm$0.014 \\
256 & 0.619 & [0.608, 0.630] & 0.0046 & 2 & $\pm$0.011 \\
512 & 0.634 & [0.623, 0.645] & 0.0052 & 2 & $\pm$0.008 \\
\bottomrule
\end{tabular}
\end{table}

\begin{table}[t]
\centering
\caption{Per-feature AUC at the 1v2 boundary by text length. All features improve monotonically, consistent with the DPI prediction that aggregate mutual information scales with sample size.}
\label{tab:perfeature_length}
\small
\begin{tabular}{lccc}
\toprule
\textbf{Feature} & \textbf{128 tok} & \textbf{256 tok} & \textbf{512 tok} \\
\midrule
freq\_entropy & 0.554 & 0.571 & 0.582 \\
surp\_std & 0.548 & 0.565 & 0.578 \\
surp\_d1\_std & 0.542 & 0.558 & 0.573 \\
hapax\_ratio & 0.536 & 0.555 & 0.569 \\
ttr & 0.533 & 0.551 & 0.566 \\
\bottomrule
\end{tabular}
\end{table}

The 128-token failure is a statistical estimation effect, not an absence of depth signal. At 128 tokens, distributional features such as hapax ratio and frequency entropy have high estimation variance ($\text{SE} \propto 1/\sqrt{n}$), collapsing the 1v2 signal into noise. The 128$\to$256 improvement ($+$2.9pp) exceeds the 256$\to$512 improvement ($+$1.5pp), indicating that 256 tokens captures the majority of statistical precision gains.

\subsection{Domain Ablation: Wikipedia}
\label{sec:domain}

We replicate the baseline experiment (Qwen2.5-1.5B base, continuation, C4) using Wikipedia seed data.

\begin{table}[t]
\centering
\caption{Domain comparison: C4 vs.\ Wikipedia. $\kstar = 2$ in both domains.}
\label{tab:domain}
\small
\begin{tabular}{lcccc}
\toprule
\textbf{Domain} & \textbf{0v1} & \textbf{1v2} & \textbf{1v2 CI lower} & $\kstar$ \\
\midrule
C4 & 0.987 & 0.626 & 0.617 & 2 \\
Wikipedia & 0.991 & 0.618 & 0.607 & 2 \\
\bottomrule
\end{tabular}
\end{table}

The 0.8pp difference in 1v2 AUC ($\Delta = -0.008$) is within confidence intervals. Feature ranking Spearman $\rho = 0.878$ between domains confirms that the degradation mechanism is domain-independent, with minor shifts in feature importance attributable to Wikipedia's more structured formatting.

\subsection{Ecological Validity: Retrain-Chain}
\label{sec:retrain}

To address the concern that prompt-chains may not represent real-world recursive contamination, we construct a retrain-chain using Pythia-1.4B with LoRA fine-tuning (rank=16, $\alpha$=32) at each depth (0--4).

\begin{table}[t]
\centering
\caption{Prompt-chain vs.\ retrain-chain on Pythia-1.4B. Both yield $\kstar = 2$. Feature ranking Spearman $\rho = 0.841$; top-3 features shared.}
\label{tab:retrain}
\small
\begin{tabular}{lcccccl}
\toprule
\textbf{Chain type} & \textbf{0v1} & \textbf{1v2} & \textbf{2v3} & \textbf{3v4} & $\kstar$ & \textbf{Feature $\rho$} \\
\midrule
Prompt-chain & 0.993 & 0.654 & 0.552 & 0.515 & 2 & -- \\
Retrain-chain & 0.985 & 0.631 & 0.558 & 0.522 & 2 & 0.841 \\
\bottomrule
\end{tabular}
\end{table}

The retrain-chain 1v2 AUC is 2.3pp lower than the prompt-chain, attributable to LoRA's low-rank constraint acting as an information-preserving regularizer that slows distributional drift. Critically, $\kstar = 2$ and top discriminative features are shared, confirming ecological validity.

\subsection{$\kstar$-Aware Data Curation}
\label{sec:curation}

We evaluate the practical utility of $\kstar = 2$ as a curation guideline by comparing three filtering strategies at contamination doses of 10\%, 30\%, and 50\%:
\begin{itemize}
  \item \textbf{Binary $\kstar$-aware}: Remove documents predicted as depth $\geq 1$ (i.e., any synthetic text).
  \item \textbf{Graduated ($\alpha = 0.9$)}: Weight documents by predicted depth using a decay schedule.
  \item \textbf{No filter}: Baseline with no contamination removal.
\end{itemize}

\begin{table}[t]
\centering
\caption{Dose-response: MMLU scores under three filtering strategies. Binary $\kstar$-aware filtering is optimal at all doses; the advantage scales with contamination level and reaches significance at 50\%.}
\label{tab:dose}
\small
\begin{tabular}{lccccc}
\toprule
\textbf{Dose} & \textbf{Binary} & \textbf{Graduated} & \textbf{No filter} & \textbf{Bin.\ vs.\ NF $\Delta$} & \textbf{Bin.\ vs.\ Grad.\ $p$} \\
\midrule
10\% & 0.5879 & 0.5876 & 0.5873 & +0.06pp & 0.42 (n.s.) \\
30\% & 0.5912 & 0.5897 & 0.5858 & +0.54pp & 0.08 (marg.) \\
50\% & 0.5941 & 0.5892 & 0.5865 & +0.76pp & 0.003 \\
\bottomrule
\end{tabular}
\end{table}

Binary filtering outperforms graduated filtering at every tested dose. At 50\% contamination, the advantage is statistically significant ($p = 0.003$, $\Delta = +0.49$pp). This directly follows from $\kstar = 2$: because depth $\geq 2$ is indistinguishable from depth 1 in the 15D feature space, graduated weighting that attempts to exploit fine-grained depth signal beyond $\kstar$ overfits to noise.

\subsection{Naturalistic Validation on Common Crawl}
\label{sec:cc}

As a suggestive (not conclusive) validation, we apply the trained OBD classifier to 50,000 documents each from Common Crawl 2025-Q1 and 2026-Q1 snapshots.

\begin{table}[t]
\centering
\caption{Predicted depth distribution in Common Crawl snapshots and correlation with independent text quality indicators.}
\label{tab:cc}
\small
\begin{tabular}{lccccc}
\toprule
\textbf{Snapshot} & \textbf{Depth 0} & \textbf{Depth 1} & \textbf{Depth 2} & \textbf{Depth 3+} & \textbf{Mean depth} \\
\midrule
CC 2025-Q1 & 72\% & 18\% & 7\% & 3\% & 0.42 \\
CC 2026-Q1 & 65\% & 22\% & 9\% & 4\% & 0.51 \\
\bottomrule
\end{tabular}
\vspace{6pt}
\begin{tabular}{lccc}
\toprule
\textbf{Quality indicator} & \textbf{Spearman $\rho$} & \textbf{$p$-value} & \textbf{Direction} \\
\midrule
AlignScore & 0.418 & $<$0.001 & Higher depth $\to$ lower quality \\
TTR & $-$0.372 & $<$0.001 & Higher depth $\to$ lower diversity \\
Hapax ratio & $-$0.341 & $<$0.001 & Higher depth $\to$ fewer rare words \\
Freq.\ entropy & $-$0.389 & $<$0.001 & Higher depth $\to$ lower entropy \\
\bottomrule
\end{tabular}
\end{table}

The predicted depth distributions are consistent with reported trends in synthetic content prevalence. The 2025$\to$2026 increase in mean predicted depth (0.42$\to$0.51, $p < 0.001$) aligns with public estimates of growing synthetic content on the web. Moderate correlations with independent quality indicators ($\rho \approx 0.34$--$0.42$) provide convergent validity.

\paragraph{Limitations.} No ground-truth depth labels exist for Common Crawl; all analysis is correlational. Non-recursive synthetic text (e.g., SEO-generated, template-based) confounds predicted depth.

% ============================================================
\section{Analysis and Discussion}
\label{sec:discussion}

\subsection{Summary of $\kstar$ Across All Conditions}

Table~\ref{tab:summary} summarizes $\kstar$ across all 19 experimental conditions. Under standard conditions ($\geq$256 tokens, $T \geq 0.9$, compatible scorer architecture), 14/14 conditions yield $\kstar = 2$.

\begin{table}[t]
\centering
\caption{Summary of $\kstar$ across all 19 tested conditions. Standard conditions (14 conditions): $\kstar = 2$ in 14/14 (100\%). Boundary conditions (5 conditions): $\kstar = 1$ in 5/5.}
\label{tab:summary}
\small
\begin{tabular}{llcc}
\toprule
\textbf{Condition category} & \textbf{Specific condition} & \textbf{$\kstar$} & \textbf{1v2 CI lower} \\
\midrule
\multirow{5}{*}{Architecture (standard)} & Pythia-1.4B & 2 & 0.645 \\
 & OLMo-1B & 2 & 0.605 \\
 & Qwen2.5-1.5B (self) & 2 & 0.617 \\
 & Qwen2.5-7B & 2 & 0.608 \\
 & LLaMA-3.1-8B & 2 & 0.607 \\
\midrule
\multirow{3}{*}{Decoding (Pythia)} & Greedy & 2 & 0.638 \\
 & Nucleus $p{=}0.9$ & 2 & 0.627 \\
 & $T{=}1.2$ & 2 & 0.654 \\
\midrule
\multirow{2}{*}{Length (Qwen-1.5B)} & 256 tokens & 2 & 0.608 \\
 & 512 tokens & 2 & 0.623 \\
\midrule
Domain & Wikipedia & 2 & 0.607 \\
\midrule
Ecological & Retrain-chain & 2 & 0.619 \\
\midrule
\multirow{2}{*}{Cross-scorer (robust gen.)} & Pythia $\times$ 4 scorers & 2 & $\geq$0.624 \\
 & OLMo $\times$ 4 scorers & 2 & $\geq$0.596\textsuperscript{*} \\
\midrule
\midrule
\multirow{5}{*}{Boundary conditions} & Mistral-7B (scorer mismatch) & 1 & 0.582 \\
 & Pythia $T{=}0.7$ (low temp.) & 1 & 0.597 \\
 & 128 tokens (short text) & 1 & 0.576 \\
 & Qwen-cont $\times$ Pythia scorer & 1 & $<$0.60 \\
 & Qwen-cont $\times$ OLMo scorer & 1 & $<$0.60 \\
\bottomrule
\end{tabular}
\vspace{2pt}
\raggedright\footnotesize \textsuperscript{*}OLMo CI lowers range 0.596--0.607; all self/Qwen-scored above $\thetaval$.
\end{table}

\subsection{Boundary Conditions of Universality}
\label{sec:boundary}

Five conditions yield $\kstar = 1$. We unify these under the DPI framework as cases where the measurable mutual information between depth and the 15D feature representation falls below the detection threshold.

\paragraph{Short text ($<$256 tokens).} At 128 tokens, distributional statistics have high estimation variance ($\text{SE} \propto 1/\sqrt{n}$). JSD at 1v2 drops 27\% from the 512-token baseline (0.0038 vs.\ 0.0052), reflecting a noise floor rather than absent signal: the 128-token CI upper bound (0.604) crosses $\thetaval$, indicating the signal exists but is buried in estimation noise. The DPI framework predicts this: per-token mutual information loss is approximately constant ($\varepsilon$), but 15D features measure \emph{aggregate} statistics whose precision requires sufficient sample size.

\paragraph{Low temperature ($T = 0.7$, Pythia).} Low temperature concentrates the output distribution over high-probability tokens, reducing the vocabulary entropy available for degradation measurement. Pythia's smaller vocabulary (50K vs.\ Qwen's 151K) makes it disproportionately sensitive: at $T = 0.7$, Pythia's effective vocabulary narrows more severely, reducing the measurable space for generational degradation features (freq\_entropy, hapax\_ratio). The Pythia--Qwen gap at $T = 0.7$ ($-$7.7pp) is nearly double the gap at other temperatures ($-$3.9 to $-$4.7pp). In DPI terms, low temperature reduces the channel capacity of each generation step, lowering the mutual information available to the feature extractor.

\paragraph{Cross-architecture scorer mismatch (Mistral).} Mistral-7B-v0.3 uses SentencePiece BPE with a sliding window attention (SWA) and grouped-query attention (GQA) architecture. When scored by Qwen2.5-1.5B (different BPE scheme, different architecture), tokenizer mismatch introduces noise into surprisal-derived features: the same text has different token boundaries under the two tokenizers, reducing surprisal feature precision. Additionally, Mistral's high-capacity architecture (GQA with 8 heads, 32K vocabulary) acts as a higher-fidelity channel with lower per-step information loss, placing its 1v2 signal closer to the detection threshold. The 1v2 AUC of 0.593 remains well above chance (0.50), indicating the signal exists but falls below the $\thetaval = 0.60$ operational threshold---a measurement precision limitation, not a physical boundary.

\paragraph{Unified interpretation.} All three failure modes reduce the same quantity: the mutual information between recursive depth and the 15D feature vector, as measured under the specific experimental conditions. The DPI framework predicts each failure \emph{a priori}: shorter sequences reduce aggregate MI through sample size; low temperature reduces MI through channel narrowing; scorer mismatch reduces MI through measurement noise. The consistency of $\kstar = 2$ when these conditions are met, and its principled degradation when they are not, provides stronger evidence for the universality of the underlying phenomenon than perfect consistency across all conditions would.

\subsection{Implications for Data Curation}

The finding that $\kstar = 2$ directly predicts the optimality of binary filtering: because depths $\geq 2$ are indistinguishable from depth 1, any graduated weighting scheme that attempts to use fine-grained depth information beyond the binary boundary is fitting noise. The dose-response experiment (\S\ref{sec:curation}) confirms this prediction across contamination levels from 10\% to 50\%, with the binary advantage scaling monotonically with dose. For practitioners, the recommendation is straightforward: use binary synthetic-vs-human filtering for data curation rather than attempting depth-aware graduated approaches.

% ============================================================
\section{Conclusion}

We present the first systematic measurement of the discrimination boundary $\kstar$ for recursive synthetic text. Across 14 standard conditions spanning six model families, four decoding strategies, two text domains, and three text lengths, $\kstar = 2$ is consistent. Three types of boundary conditions---short texts ($<$256 tokens), low temperature ($T < 0.9$), and cross-architecture scorer mismatch---predictably degrade $\kstar$ to 1, each explained by the DPI framework as a reduction in measurable mutual information.

These results have both theoretical and practical implications. Theoretically, the universality of $\kstar = 2$ under standard conditions suggests a fundamental property of iterative text generation that transcends specific model architectures. Practically, $\kstar = 2$ provides an actionable guideline: binary contamination filtering at the $\kstar$ boundary is the optimal curation strategy, eliminating the need for complex graduated depth-weighting schemes.

\paragraph{Limitations.} Our experiments cover models from 1B to 8B parameters; scaling behavior at 13B+ remains untested. The 15D feature set, while model-free and transferable, may underestimate $\kstar_{\text{theory}}$---neural probe experiments (not reported here for space) suggest $\kstar \geq 3$ for model-specific representations, at the cost of universality. The Common Crawl validation is correlational without ground-truth depth labels. A continuous dose-response curve (beyond three dose points) is left to future work.

\paragraph{Future work.} Extensions to larger models ($\geq$13B), code and multilingual domains, and richer feature representations (including learned, model-specific features) would further characterize the discrimination boundary. Integration with existing data pipelines and evaluation on diverse downstream tasks beyond MMLU would strengthen the practical impact.

\bibliographystyle{plainnat}
\bibliography{references}

\end{document}
